\subsection*{Finite Truth Table Intuition}

\begin{remark*}[Finite truth table intuition]
Truth tables are finite because a single formula contains only finitely many
variables. More generally, a finite set of formulas contains only finitely many
variables, so its satisfiability can be checked by a finite truth table over
those variables.
\end{remark*}

\begin{remark*}[Finite test]
For a finite \(\Delta\subseteq\WFF\), the variables appearing in \(\Delta\) are
finite in number. A truth table over those variables decides whether
\(\Delta\) is satisfiable.
\end{remark*}

\begin{remark*}[Why compactness is surprising]
Finite truth tables directly decide finite sets. Compactness says that, for
propositional logic, passing all finite tests is enough to pass the entire
possibly infinite test.
\end{remark*}

\begin{dependencies}
\begin{itemize}
  \item \hyperref[lem:finiteness-of-variables-propositional-formulas]{Finiteness of Variables}
  \item \hyperref[thm:finite-truth-table-propositional-logic]{Finite Truth Table}
  \item \hyperref[def:finitely-satisfiable-set-of-propositional-formulas]{Finitely Satisfiable Set of Propositional Formulas}
\end{itemize}
\end{dependencies}
